En este capítulo se resume la estrategia de validación utilizada en la iteración actual del proyecto. Dado que se trata de un \textit{TFG} con una aplicación \textit{full-stack} de alcance incremental, la verificación combina comprobaciones automáticas de regresión con validaciones funcionales manuales y asistidas sobre los flujos críticos de los módulos \texttt{M1}, \texttt{M2}, \texttt{M3} y \texttt{M4}, y sobre los bloques operativos ya disponibles de \texttt{M5} y \texttt{M6}.

\section{Estrategia}

La estrategia de pruebas seguida persigue dos objetivos complementarios:

\begin{itemize}
    \item detectar regresiones técnicas rápidas después de cada cambio relevante mediante comprobaciones automáticas;
    \item confirmar de forma funcional que los recorridos visibles para cada rol continuan funcionando de acuerdo con el alcance operativo implementado.
\end{itemize}

Las regresiones técnicas se evaluan repitiendo siempre la misma secuencia base:

\begin{enumerate}
    \item comprobación estática de tipos con \texttt{TypeScript};
    \item análisis estático de código con \texttt{ESLint};
    \item construcción completa de producción con \texttt{Next.js};
    \item smoke test funcional de \texttt{M4} sobre servicios y persistencia;
    \item pasada visual hidratada de \texttt{M4} sobre rutas críticas por rol;
    \item comprobación puntual de \texttt{M5} creando una ficha técnica temporal, registrando un servicio con tonalidad e imágenes finales, generando su borrador técnico de correo, creando una plantilla reutilizable y limpiando el dato generado;
    \item comprobación puntual de \texttt{M6} validando migración, carga de pantallas de comunicación, visibilidad de promociones activas, aplicación de descuentos promocionales en pedidos, notificaciones segmentadas, recordatorios e inscripción a formaciones publicadas.
\end{enumerate}

Esta rutina ofrece una señal rápida sobre errores de tipado, problemas de compilación, importaciones rotas, uso incorrecto de componentes de servidor o cliente, fallos en rutas dinámicas y desviaciones funcionales del flujo de pedidos, reparto y cobro.

\section{Pruebas Unitarias}

\textcolor{red}{En la versión actual no se ha incorporado todavía una batería formal de pruebas unitarias con un \textit{runner} dedicado.} Aún así, parte de la validación de bajo nivel se apoya en tres mecanismos que reducen errores locales antes de integrar cambios:

\begin{itemize}
    \item el tipado estricto de \texttt{TypeScript} sobre contratos, componentes y servicios;
    \item la centralizacion de reglas de negocio en servicios de dominio de usuarios, operativa comercial, catálogo y pedidos;
    \item las validaciones normalizadas de entrada antes de persistir cambios, lo que limita estados inconsistentes en operaciones de administración, catálogo, reparto, cobro y seguimiento técnico del salón.
\end{itemize}

\textcolor{red}{La ausencia de una suite unitaria completa sigue siendo una mejora pendiente}, pero la separación actual de servicios facilita su incorporación futura.

\section{Pruebas de Integración}

Las pruebas de integración se centran en comprobar la interacción entre interfaz, \textit{Route Handlers}, autenticación, persistencia y servicios externos. En esta línea resultan especialmente relevantes los siguientes puntos:

\begin{itemize}
    \item autenticación con \texttt{Auth.js} y control de acceso por rol en \texttt{proxy.ts};
    \item coordinación entre formularios administrativos y servicios \texttt{TypeORM} para usuarios, clientes, visitas, catálogo, pedidos, seguimiento técnico del salón y comunicaciones internas;
    \item coordinación entre el centro administrativo de comunicaciones, las APIs de \texttt{M6}, los servicios de pedido y las vistas cliente/comercial de promociones, formaciones, avisos y recordatorios;
    \item integración con \texttt{Nominatim} y \texttt{OSRM} para geocodificación y cálculo de rutas del módulo \texttt{M2};
    \item integración con \texttt{Cloudinary} para subida, reemplazo y validación de imágenes de perfil y catálogo en \texttt{M1} y \texttt{M3};
    \item integración con \texttt{QuickChart} para representar el QR operativo de cada pedido dentro del flujo de reparto de \texttt{M4}.
\end{itemize}

Estas comprobaciones se realizan principalmente a través de la propia aplicación en entorno local, verificando que una acción iniciada desde la interfaz produce el efecto esperado en la API y en la base de datos.

\section{Pruebas de Sistema}

La evidencia automatizada disponible para el estado actual del repositorio se ha obtenido ejecutando los siguientes comandos sobre la aplicación \texttt{app-tfg/}:

\begin{itemize}
    \item \texttt{npm run typecheck}
    \item \texttt{npm run lint}
    \item \texttt{npm run build}
    \item \texttt{npm run m1-m3:closeout}
    \item \texttt{npm run m4:closeout}
    \item \texttt{npm run m4:ui-check}
    \item \texttt{npm run m5:technical-email-smoke}
    \item \texttt{npm run m5:service-templates-smoke}
    \item \texttt{npm run m5:salon-visual-smoke}
    \item \texttt{npm run m6:communications-smoke}
    \item \texttt{npm run m6:promotion-discount-smoke}
    \item \texttt{npm run m6:client-tier-smoke}
    \item \texttt{npm run m6:closeout}
    \item \texttt{npm run m7:integrations-smoke}
    \item \texttt{npm run m7:support-smoke}
    \item \texttt{npm run m7:operations-smoke}
    \item \texttt{npm run m7:audit-export-smoke}
    \item \texttt{npm run m7:rate-limit-smoke}
    \item \texttt{npm run m7:operations-report-smoke}
    \item \texttt{npm run m7:closeout}
    \item \texttt{npm run migration:run}
\end{itemize}

El resultado de esta validación es satisfactorio: la comprobación de tipos, el análisis estático, la construcción completa de producción y la aplicación de migraciones finalizan sin errores. Además, se ha incorporado \texttt{npm run m1-m3:closeout}, una comprobación específica para cubrir el hueco de evidencia técnica de los módulos iniciales. Este script crea datos temporales y los elimina al finalizar, validando \texttt{M1}, \texttt{M2} y \texttt{M3} con \texttt{11/11} comprobaciones satisfactorias: alta directa de comercial, solicitud pública de alta, aprobación de cliente, autenticación por credenciales, cambio de contraseña, ficha de cliente con geodatos, asignación comercial-cliente, visita, ruta y estructura mínima de catálogo con categoría, línea, subcategoría jerárquica y producto.

El cierre técnico de \texttt{M4} queda reforzado por dos evidencias adicionales:

\begin{itemize}
    \item un smoke test funcional de servicios y persistencia, que cubre el flujo principal de pedidos, reparto y cobro y finaliza con \texttt{19/19} comprobaciones satisfactorias en la batería principal y \texttt{12/12} en la pasada de cierre;
    \item una verificación visual hidratada sobre cliente, comercial y administrador, reutilizable para comprobar bandeja de reparto, detalle de visita, cobros, ETA del cliente, detalle administrativo y gestión de errores de sesión o de carga.
\end{itemize}

Como validación manual adicional de \texttt{M5}, se ha ejecutado una pasada completa en navegador local sobre \texttt{/clients/salon-clients}: autenticación como cliente, alta de ficha técnica, guardado y aplicación de plantilla con tonalidad \texttt{BLUE GUMMY}, subida real de imagen final a \texttt{Cloudinary}, registro de servicio, filtros de historial, generación de correo técnico, edición del servicio, retirada de imagen, borrado del servicio, borrado de plantilla, comprobación responsive móvil y control de acceso sin sesión, con administrador, con comercial y con un segundo cliente temporal. La pasada no registra errores de consola, excepciones de ejecución, fallos de red bloqueantes ni errores HTTP inesperados.

Una mejora técnica asociada a esta fase de pruebas ha consistido en desacoplar la construcción de \texttt{Next.js} de fuentes remotas, de modo que \texttt{npm run build} ya no depende de descargar \texttt{Geist} desde \texttt{next/font/google} para completarse en local.

Como comprobación transversal adicional de cierre, se ha ejecutado un barrido ampliado en navegador headless sobre 68 rutas, pantallas dinámicas, redirecciones por rol y consultas de API de los módulos \texttt{M1} a \texttt{M5}: autenticación, registro público, sesión, redirección de rutas protegidas, restricciones de acceso por rol, perfil, cambio de contraseña, panel de cliente, administración de clientes y asignaciones, detalle de cliente, operativa comercial, visitas, rutas, catálogo, formularios administrativos, coloración, detalles de productos y cartas de color, pedidos, cobros y fichas técnicas del salón. El resultado ha sido \texttt{68/68} comprobaciones satisfactorias, desglosadas en \texttt{M1 12/12}, \texttt{M2 19/19}, \texttt{M3 24/24}, \texttt{M4 9/9} y \texttt{M5 4/4}, sin errores de consola, excepciones de ejecución, fallos de red bloqueantes ni errores HTTP inesperados.

Para \texttt{M6}, la validación de esta iteración ha incluido la aplicación correcta de la migración \texttt{039}, la carga de las nuevas rutas de comunicación en entorno de producción local, la visualización de una promoción activa temporal desde los roles cliente y comercial, la aplicación de descuentos porcentuales activos sobre borradores y pedidos cuando coinciden cliente, segmento, producto o línea comercial, la visualización del rango comercial del cliente en su perfil, la consulta de formaciones publicadas, la revisión de avisos internos pendientes y la inscripción real a una formación publicada temporal. Además, se ha incorporado \texttt{npm run m6:communications-smoke}, que crea un segmento temporal, edita sus datos, asigna un cliente, publica una promoción segmentada, comprueba visibilidad por rol, valida que las notificaciones internas respetan el segmento, archiva la promoción, publica una formación, registra y cancela una inscripción, crea y actualiza un recordatorio y elimina todos los datos temporales generados. El cierre queda reforzado con \texttt{npm run m6:promotion-discount-smoke}, que crea una promoción temporal de descuento vinculada a cliente y producto, comprueba que el descuento se aplica en el borrador comercial y que el pedido final conserva el porcentaje, el total de línea y el total del pedido, eliminando después la promoción y el pedido generados. La comprobación \texttt{npm run m6:client-tier-smoke} valida que todos los clientes tienen un rango comercial base, que el servicio del badge reconoce el rango \textit{Platino}, conserva los segmentos asignados y expone ventajas comerciales sin dejar asignaciones temporales en base de datos. El script \texttt{npm run m6:closeout} encadena las tres comprobaciones para repetir el cierre completo del módulo antes de iniciar nuevas iteraciones. Como comprobación visual adicional, se ha iniciado sesión con usuarios reales de prueba de los tres perfiles y se han cargado \path{/admin/communications}, \path{/clients}, \path{/profile}, \path{/clients/promotions}, \path{/clients/trainings}, \path{/clients/notifications}, \path{/commercials/promotions} y \path{/commercials/trainings} en producción local. La revisión de permisos ha comprobado que las APIs administrativas devuelven \texttt{401} sin sesión o con usuario no administrador, que el administrador accede correctamente al centro de segmentos con \texttt{200}, que el administrador queda bloqueado en las APIs de comunicación de usuario con \texttt{403}, y que cliente y comercial acceden a promociones, formaciones y notificaciones con \texttt{200}. También se ha validado desde la interfaz el borrado con confirmación inline de una promoción, una formación y un segmento temporales, verificando posteriormente que la base de datos no conservaba esos registros. No se observaron errores de consola ni fallos de carga en las páginas comprobadas. La creación visual de recordatorios queda cubierta por tipado, API, build y smoke de servicio, aunque durante la comprobación manual el navegador integrado no permitió escribir en determinados formularios por una limitación de su clipboard virtual.

Para \texttt{M7}, el primer incremento de esta iteración incorpora el inventario técnico de integraciones de soporte. La comprobación \texttt{npm run m7:integrations-smoke} valida que el inventario contiene las integraciones esperadas, que no existen identificadores duplicados, que cada elemento expone estado, configuración, uso operativo y degradación prevista, y que el resumen agregado de estados es coherente. Este inventario se conserva como apoyo de desarrollo y documentación técnica, no como pantalla ordinaria de administración.

El segundo incremento de \texttt{M7} incorpora el inventario de soporte técnico y compatibilidad. La comprobación \texttt{npm run m7:support-smoke} valida que el inventario contiene las capacidades esperadas de navegador compatible, manifiesto PWA, \textit{service worker} y formularios seguros ante hidratación; también comprueba que cada capacidad expone estado, evidencias operativas y degradación prevista.

El tercer incremento de \texttt{M7} incorpora el resumen técnico de operación transversal. La comprobación \texttt{npm run m7:operations-smoke} valida que el resumen contiene los bloques de auditoría, rate limiting, integraciones, soporte técnico y operación empresarial, y que las métricas agregadas son coherentes. Tras la revisión final, esta información queda fuera de la navegación principal para evitar que el administrador de negocio tenga que consultar pantallas de diagnóstico.

El cuarto incremento de \texttt{M7} refuerza la trazabilidad administrativa mediante exportación \texttt{CSV} desde \path{/admin/audit}. La comprobación \texttt{npm run m7:audit-export-smoke} valida que se generan los ficheros de accesos y acciones administrativas con cabeceras, tipo de contenido y nombres de archivo consistentes.

El quinto incremento de \texttt{M7} incorpora el diagnóstico agregado de rate limiting como infraestructura técnica. La comprobación \texttt{npm run m7:rate-limit-smoke} valida que se diagnostican las políticas esperadas, que cada una expone límites efectivos, valores por defecto y estado operativo, y que el resumen agregado es coherente. La configuración editable de estas políticas no se expone al administrador ordinario; los endpoints asociados quedan desactivados salvo activación expresa mediante entorno técnico.

El sexto incremento de \texttt{M7} añade la generación técnica del informe operativo. La comprobación \texttt{npm run m7:operations-report-smoke} valida que el informe contiene la cabecera esperada, que incluye los cinco bloques transversales y que conserva un número suficiente de métricas para cada sección.

El séptimo incremento de \texttt{M7} materializa el registro empresarial persistente previsto en el análisis. La comprobación \texttt{npm run m7:enterprise-operations-smoke} valida que existen configuraciones e integraciones persistidas, que la integración \textit{Factusol mediante n8n} queda preparada, que se puede registrar una operación de integración y que se puede generar una propuesta de pedido a proveedor con líneas calculadas a partir de productos activos. El script limpia los datos temporales generados y restaura el estado inicial de la integración probada.

El script \texttt{npm run m7:closeout} encadena las comprobaciones de integraciones, soporte, operación, exportación de auditoría, rate limiting, informe operativo y operación empresarial para repetir el cierre del módulo completo.

Como cierre funcional de \texttt{M7}, se ha ejecutado una pasada completa de validación con \texttt{npm run m7:closeout}, \texttt{npm run typecheck}, \texttt{npm run lint}, \texttt{npm run build}, \texttt{npm run migration:run} y la compilación \LaTeX{} de la memoria. La revisión manual en producción local se centra en las pantallas finales que sí quedan en el flujo de administración: \path{/admin/audit} y \path{/admin/settings}. Las rutas de soporte, integraciones y operación transversal se mantienen redirigidas fuera de la navegación ordinaria, mientras que sus servicios siguen cubiertos por scripts de validación. Además, se han comprobado los endpoints administrativos principales sin sesión, confirmando respuesta \texttt{401}, y no se han observado errores de consola en las pantallas revisadas.

Como comprobación no funcional adicional, se ha ejecutado una auditoría de accesibilidad y legibilidad con WAVE Evaluation Tool \texttt{3.3.1.0} sobre la aplicación levantada en producción local mediante \texttt{npm run build} y \texttt{npm run start -- --port 3000}. Dado que la API remota de WAVE requiere clave de uso y una URL accesible externamente, para auditar el entorno local se ha descargado la extensión oficial de WAVE, se ha extraído su motor de evaluación y se ha ejecutado sobre Chrome headless mediante el protocolo de depuración del navegador. La pasada ha cubierto login público, panel de administración, operaciones, rutas comerciales, pedidos, panel de cliente, fichas de salón y panel de cliente en viewport móvil.

Antes de la pasada final se corrigieron los problemas detectados: asociación explícita de etiquetas \texttt{label} en login y fichas de salón, incorporación de \texttt{legend} accesible en \texttt{SafeForm}, nombres accesibles y estado activo en la nueva barra lateral por rol, ocultación real de la barra lateral móvil cuando está cerrada, preservación acotada de la transición del menú lateral en navegadores con \texttt{prefers-reduced-motion}, eliminación de títulos redundantes, ajuste del texto alternativo de avatares y aumento de textos de \texttt{10px} en cabeceras, tarjetas y badges.

\begin{table}[H]
    \centering
    \small
    \renewcommand{\arraystretch}{1.15}
    \begin{tabularx}{\linewidth}{|X|p{0.23\linewidth}|c|c|c|}
        \hline
        \textbf{Vista auditada} & \textbf{Ruta} & \textbf{Errores} & \textbf{Contraste} & \textbf{Avisos} \\
        \hline
        Login público & \path{/login} & 0 & 0 & 0 \\
        \hline
        Panel admin & \path{/admin} & 0 & 0 & 0 \\
        \hline
        Auditoría admin & \path{/admin/audit} & 0 & 0 & 4 \\
        \hline
        Rutas comercial & \path{/commercials/routes} & 0 & 0 & 0 \\
        \hline
        Pedidos comercial & \path{/commercials/orders} & 0 & 0 & 4 \\
        \hline
        Panel cliente & \path{/clients} & 0 & 0 & 1 \\
        \hline
        Fichas de salón & \path{/clients/salon-clients} & 0 & 0 & 0 \\
        \hline
        Panel cliente móvil & \path{/clients} & 0 & 0 & 1 \\
        \hline
    \end{tabularx}
    \caption{Resultado de auditoría WAVE sobre rutas representativas}
    \label{tab:wave-accessibility-audit}
\end{table}

Los avisos residuales de WAVE corresponden a \textit{Possible heading} en valores destacados de tarjetas de métricas. Se han revisado manualmente y se mantienen como valores visuales, no como encabezados de sección, por lo que no se consideran bloqueantes para el estado actual. No se registran errores de accesibilidad, errores de contraste ni avisos de texto demasiado pequeño en la pasada final.

\subsection{Pruebas Funcionales}

Las comprobaciones funcionales prioritarias se organizan por módulo:

\begin{itemize}
    \item \texttt{M1}: registro de solicitud de alta, autenticación con credenciales, navegación protegida por rol, edición de perfil y cambio de contraseña.
    \item \texttt{M2}: administración de clientes, asignación comercial-cliente, geolocalización confirmada, creación y seguimiento de visitas, aplazamientos automáticos con aviso al comercial, avisos al cliente por creación o reubicación, configuración diaria del comercial y cálculo de ETA para cliente.
    \item \texttt{M3}: mantenimiento del catálogo desde administración, consulta pública de productos activos, filtrado por categorías y líneas, acceso a fichas técnicas, consulta de recursos de apoyo y exploración de cartas de color con sus referencias.
    \item \texttt{M4}: borradores de pedido en cliente y comercial, confirmación y cancelación coherente, bandeja de pedidos pendientes de reparto, visita de entrega con validación QR, paso a \texttt{delivered}, subflujo mínimo de cobro y trazabilidad cruzada desde admin.
    \item \texttt{M5}: alta de clientes del salón, edición de ficha base, registro, edición y borrado de servicios con ficha técnica, selección de tonalidades de color, galería de imágenes finales, plantillas reutilizables, filtrado del historial, uso de productos, sugerencias derivadas y generación de borradores técnicos editables con persistencia real sobre la base de datos.
    \item \texttt{M6}: gestión, edición y borrado de promociones, segmentos y formaciones desde administración, visualización de rango comercial del cliente, visibilidad de promociones activas por rol, aplicación de descuentos promocionales sobre borradores y pedidos, inscripción a formaciones publicadas, notificaciones internas segmentadas, correo SMTP, Web Push PWA y recordatorios personales.
    \item \texttt{M7}: consulta administrativa de auditoría, exportación \texttt{CSV} de accesos y acciones administrativas, configuración visible de canales de avisos automáticos, diagnóstico técnico de rate limiting, inventarios internos de integraciones y soporte, informe operativo técnico, registro persistente de integraciones empresariales, operaciones de intercambio y propuestas de pedido a proveedor.
\end{itemize}

En \texttt{M4}, la regresión funcional más sensible consiste en comprobar que el mismo pedido se mantiene coherente al pasar por tres contextos distintos: cliente, comercial y administrador. Por ello, cualquier cambio en entidades, serializadores, servicios de pedido o lógica de visitas de reparto debe revalidarse en esos tres recorridos.

\subsection{Trazabilidad entre módulos y evidencias}

La Tabla~\ref{tab:test-traceability} relaciona requisitos funcionales representativos con elementos verificables de la implementación. La finalidad de esta tabla es dejar constancia de que cada flujo documentado se apoya en pantallas, rutas API, entidades persistentes y pruebas ejecutables, evitando que la memoria quede únicamente como una descripción narrativa del sistema.

\begin{landscape}
{\scriptsize
\setlength{\tabcolsep}{1pt}
\renewcommand{\arraystretch}{1.18}
\begin{longtable}{|>{\RaggedRight\arraybackslash}p{0.075\linewidth}|>{\RaggedRight\arraybackslash}p{0.055\linewidth}|>{\RaggedRight\arraybackslash}p{0.145\linewidth}|>{\RaggedRight\arraybackslash}p{0.165\linewidth}|>{\RaggedRight\arraybackslash}p{0.185\linewidth}|>{\RaggedRight\arraybackslash}p{0.180\linewidth}|>{\RaggedRight\arraybackslash}p{0.155\linewidth}|}
\caption{Trazabilidad técnica entre requisitos, implementación y evidencias de prueba}
\label{tab:test-traceability}\\
\hline
\textbf{Requisito} & \textbf{Módulo} & \textbf{Ruta / pantalla} & \textbf{Endpoint} & \textbf{Entidad de base de datos} & \textbf{Prueba realizada} & \textbf{Evidencia / captura} \\
\hline
\endfirsthead
\hline
\textbf{Requisito} & \textbf{Módulo} & \textbf{Ruta / pantalla} & \textbf{Endpoint} & \textbf{Entidad de base de datos} & \textbf{Prueba realizada} & \textbf{Evidencia / captura} \\
\hline
\endhead
\texttt{RF-M1-01} & \texttt{M1} & \path{/register}, \path{/admin/users/requests/[id]}, \path{/admin/users/register} & \path{/api/auth/register-request}, \path{/api/admin/user-requests/[id]/approve}, \path{/api/admin/register-user} & \texttt{user\_requests}, \texttt{users}, \texttt{clients}, \texttt{commercials}, \texttt{client\_commercial\_assignments}, \texttt{user\_management\_log} & \texttt{npm run m1-m3:closeout}: crea comercial temporal, registra solicitud pública de cliente y la aprueba con asignación comercial. & Salida \texttt{OK M1-M3 closeout smoke: 11/11}. \\
\hline
\texttt{RF-M1-02}, \texttt{RF-M1-06} & \texttt{M1} & \path{/login}, \path{/profile}, \path{/profile/change-password} & \path{/api/auth/[...nextauth]}, \path{/api/profile}, \path{/api/account/change-password} & \texttt{users}, \texttt{user\_access\_log}, \texttt{user\_management\_log} & \texttt{npm run m1-m3:closeout}: valida usuario activo con \texttt{findUserForLogin}, contraseña cifrada con \texttt{bcrypt} y cambio de contraseña persistido. & Salida \texttt{PASS M1 autenticación local por credenciales validada} y \texttt{PASS M1 cambio de contraseña persistido}. \\
\hline
\texttt{RF-M2-01}, \texttt{RF-M2-02} & \texttt{M2} & \path{/admin/clients}, \path{/admin/clients/assignments}, \path{/commercials/clients/[id]} & \path{/api/admin/clients}, \path{/api/admin/client-commercial-assignments}, \path{/api/commercial/clients/[id]}, \path{/api/clients/[id]} & \texttt{clients}, \texttt{commercials}, \texttt{client\_commercial\_assignments}, \texttt{client\_customer\_segments} & \texttt{npm run m1-m3:closeout}: actualiza ficha de cliente con geodatos, franja horaria y asignación comercial activa. & Salida \texttt{PASS M2 ficha de cliente con geodatos y franja horaria} y \texttt{PASS M2 asignación comercial-cliente consultable}. \\
\hline
\texttt{RF-M2-03}, \texttt{RF-M2-04}, \texttt{RF-M2-05} & \texttt{M2} & \path{/commercials/visits}, \path{/commercials/visits/[id]}, \path{/commercials/routes}, \path{/commercials/settings} & \path{/api/commercial/visits}, \path{/api/commercial/visits/[id]}, \path{/api/commercial/routes/preview}, \path{/api/admin/commercial-routes/add-visit}, \path{/api/commercial/profile} & \texttt{commercial\_visits}, \texttt{commercial\_routes}, \texttt{route\_visits}, \texttt{commercials}, \texttt{app\_notifications} & \texttt{npm run m1-m3:closeout}: configura jornada del comercial, crea visita rutinaria, consulta por cliente y comercial y la incorpora a una ruta ordenada. & Salida \texttt{PASS M2 visita comercial visible por comercial y cliente} y \texttt{PASS M2 ruta comercial con visita ordenada}; comprobación visual de visitas y rutas. \\
\hline
\texttt{RF-M3-01}, \texttt{RF-M3-02}, \texttt{RF-M3-04} & \texttt{M3} & \path{/admin/catalog}, \path{/admin/catalog/products}, \path{/clients/catalog}, \path{/commercials/catalog} & \path{/api/admin/catalog/product-categories}, \path{/api/admin/catalog/product-lines}, \path{/api/admin/catalog/product-subcategories}, \path{/api/admin/catalog/products} & \texttt{product\_categories}, \texttt{product\_lines}, \texttt{product\_subcategories}, \texttt{products}, \texttt{product\_statuses} & \texttt{npm run m1-m3:closeout}: crea categoría, línea, subcategoría padre, subcategoría hija y producto activo con relaciones recargadas. & Salida \texttt{PASS M3 catálogo con categoría, línea, jerarquía y producto}; comprobación visual de catálogo y coloración. \\
\hline
\texttt{RF-M3-03} & \texttt{M3} & \path{/clients/coloration}, \path{/clients/coloration/[colorChartId]}, \path{/commercials/coloration} & \path{/api/admin/catalog/color-charts}, \path{/api/admin/catalog/color-references}, lecturas de catálogo en cliente y comercial & \texttt{color\_charts}, \texttt{color\_references}, \texttt{products} & Barrido visual hidratado y build de producción sobre rutas de coloración y detalle de carta. & Validación headless ampliada \texttt{68/68}; captura \path{client-coloration.png}. \\
\hline
\texttt{RF-M4-01}, \texttt{RF-M4-04} & \texttt{M4} & \path{/clients/orders}, \path{/clients/orders/[id]}, \path{/commercials/orders}, \path{/admin/orders/[id]} & \path{/api/clients/orders}, \path{/api/clients/orders/[id]}, \path{/api/commercial/orders}, \path{/api/admin/orders/[id]} & \texttt{orders}, \texttt{order\_lines}, \texttt{order\_statuses}, \texttt{products}, \texttt{color\_references} & \texttt{npm run m4:closeout}: crea pedido desde comercial, valida estados y consulta posterior por cliente. & Salida \texttt{OK M4 closeout smoke: 12/12}; captura \path{client-orders.png}. \\
\hline
\texttt{RF-M4-02} & \texttt{M4} & \path{/commercials/visits/[id]}, \path{/commercials/routes}, \path{/clients} & \path{/api/commercial/visits/[id]}, \path{/api/clients/delivery-estimate} & \texttt{commercial\_visits}, \texttt{orders}, \texttt{app\_notifications} & \texttt{npm run m4:closeout}: vincula pedido confirmado a reparto, valida ETA visible, aplazamiento, cancelación y QR de entrega. & Salida \texttt{PASS ETA visible con reparto real vinculado}; capturas \path{commercial-visits.png} y \path{commercial-routes-map.png}. \\
\hline
\texttt{RF-M4-03} & \texttt{M4} & \path{/commercials/orders/[id]}, \path{/admin/orders/[id]} & \path{/api/commercial/orders/[id]}, \path{/api/admin/orders/[id]} & \texttt{orders}, \texttt{order\_payment\_statuses} & \texttt{npm run m4:ui-check} y \texttt{npm run m4:closeout}: comprueba pedidos entregados, estado mínimo de cobro y trazabilidad cruzada. & Evidencia \texttt{m4-ui-evidence}; build de producción. \\
\hline
\texttt{RF-M5-01}, \texttt{RF-M5-02}, \texttt{RF-M5-07} & \texttt{M5} & \path{/clients/salon-clients}, \path{/clients/salon-clients/[id]} & \path{/api/clients/salon-clients}, \path{/api/clients/salon-clients/[id]/services}, \path{/api/clients/salon-service-templates}, \path{/api/clients/salon-service-result-images} & \texttt{salon\_clients}, \texttt{salon\_services}, \texttt{salon\_service\_templates}, \texttt{salon\_service\_result\_images}, \texttt{salon\_service\_product\_usage} & \texttt{m5:service-templates-smoke} y \texttt{m5:salon-visual-smoke}: crea ficha, servicio, plantilla e imagen final temporal. & Salidas \texttt{M5 service templates smoke OK} y \texttt{M5 salón visual smoke OK}. \\
\hline
\texttt{RF-M5-03}, \texttt{RF-M5-08} & \texttt{M5} & \path{/clients/salon-clients/[id]} & \path{/api/clients/salon-clients/[id]/services/[serviceId]/technical-email} & \texttt{salon\_services}, \texttt{salon\_service\_technical\_sheets}, \texttt{color\_references} & \texttt{m5:technical-email-smoke}: registra ficha técnica con tonalidad y genera borrador técnico de correo. & Salida \texttt{M5 technical email smoke OK}. \\
\hline
\texttt{RF-M6-01}, \texttt{RF-M6-02} & \texttt{M6} & \path{/admin/communications}, \path{/clients/promotions}, \path{/commercials/promotions} & \path{/api/admin/communications/promotions}, \path{/api/admin/communications/segments}, \path{/api/communications/promotions} & \texttt{promotions}, \texttt{customer\_segments}, \texttt{client\_customer\_segments}, \texttt{app\_notifications} & \texttt{m6:communications-smoke} y \texttt{m6:promotion-discount-smoke}: promoción segmentada, visibilidad por rol y descuento aplicado a pedido. & Salidas \texttt{M6 communications smoke OK}, \texttt{M6 promotion discount smoke OK}; comprobación visual de comunicaciones y descuento. \\
\hline
\texttt{RF-M6-03}, \texttt{RF-M6-04} & \texttt{M6} & \path{/clients/notifications}, \path{/clients/trainings}, \path{/commercials/trainings} & \path{/api/communications/notifications}, \path{/api/communications/reminders}, \path{/api/communications/trainings}, \path{/api/communications/trainings/[id]/enrollment} & \texttt{app\_notifications}, \texttt{app\_reminders}, \texttt{training\_events}, \texttt{training\_enrollments} & \texttt{m6:communications-smoke}: publica formación, registra/cancela inscripción, crea recordatorio y limpia datos temporales. & Salida \texttt{M6 communications smoke OK}; comprobación visual de avisos y formaciones. \\
\hline
\texttt{RF-M7-01}, \texttt{RF-M7-06} & \texttt{M7} & \path{/admin/audit}, \path{/admin/settings} & \path{/api/admin/audit/export}, \path{/api/admin/settings/notifications} & \texttt{user\_access\_log}, \texttt{user\_management\_log}, configuración de canales de avisos & \texttt{m7:audit-export-smoke}: exporta \texttt{CSV}; revisión funcional de ajustes de avisos. & Salida \texttt{M7 audit export smoke OK}. \\
\hline
\texttt{RF-M7-03}, \texttt{RF-M7-04}, \texttt{RF-M7-05} & \texttt{M7} & Servicios internos de integración & \path{/api/admin/integrations/status}, \path{/api/admin/enterprise-operations}, \path{/api/admin/enterprise-operations/operations}, \path{/api/admin/enterprise-operations/supplier-proposals} & \texttt{system\_configurations}, \texttt{external\_integrations}, \texttt{integration\_operations}, \texttt{supplier\_order\_proposals}, \texttt{supplier\_order\_proposal\_lines} & \texttt{m7:integrations-smoke} y \texttt{m7:enterprise-operations-smoke}: inventario técnico, operación de integración y propuesta a proveedor. & Salidas \texttt{M7 integrations smoke OK} y \texttt{M7 enterprise operations smoke OK}. \\
\hline
\texttt{RF-M7-07}, \texttt{RF-M7-08} & \texttt{M7} & \path{/admin/settings}, \path{/unsupported-browser} & \path{/api/admin/settings/rate-limits/status}, \path{/api/admin/support/status} & \texttt{app\_rate\_limit\_policies}; manifiesto PWA y \textit{service worker} sin entidad relacional propia & \texttt{m7:rate-limit-smoke} y \texttt{m7:support-smoke}: diagnostica infraestructura técnica de rate limiting, manifiesto, \textit{service worker} y compatibilidad. & Salidas \texttt{M7 rate limit smoke OK}, \texttt{M7 support smoke OK}; Tabla~\ref{tab:wave-accessibility-audit}. \\
\hline
\end{longtable}
}
\end{landscape}

\subsection{Pruebas No Funcionales}

En el estado actual del proyecto, las comprobaciones no funcionales se concentran en:

\begin{itemize}
    \item seguridad básica de acceso mediante autenticación y autorización por roles;
    \item robustez de compilación y tipado de toda la aplicación;
    \item validación básica de errores de sesión y de carga en rutas críticas de \texttt{M4};
    \item compatibilidad mínima de navegador a través de la ruta informativa para clientes no soportados;
    \item validación de tamaño, tipo y procedencia de imágenes subidas a \texttt{Cloudinary}.
\end{itemize}

\textcolor{red}{No se ha abordado todavía una campaña específica de rendimiento, carga o pruebas automatizadas de interfaz con un \textit{framework} dedicado, por lo que estas áreas permanecen como trabajo de mejora futura.}

\section{Pruebas de Aceptación}

La aceptación de la versión actual del sistema se interpreta como la disponibilidad integrada y estable de los módulos \texttt{M1}, \texttt{M2}, \texttt{M3} y \texttt{M4}, junto con bloques parciales pero operativos de \texttt{M5} y \texttt{M6}. En términos prácticos, esto implica que:

\begin{itemize}
    \item el administrador puede gestionar usuarios, clientes, catálogo, pedidos globales y el centro interno de comunicaciones;
    \item el comercial puede consultar su operativa diaria, el catálogo, sus pedidos, los repartos asociados, el subflujo mínimo de cobro, promociones, formaciones y avisos;
    \item el cliente puede consultar su previsión de reparto, su perfil, el catálogo, la coloración, el historial de pedidos, promociones, formaciones y avisos;
    \item el cliente profesional puede además mantener fichas técnicas de salón, registrar servicios con tonos e imágenes finales, reutilizar plantillas técnicas y preparar borradores técnicos editables por servicio;
    \item el flujo \texttt{pedido -> reparto -> entrega -> cobro mínimo} queda cerrado y trazable en la implementación actual;
    \item el núcleo de comunicación interna \texttt{promoción -> publicación -> consulta por rol}, \texttt{promoción -> pedido -> descuento aplicado} y \texttt{formación -> publicación -> inscripción} queda operativo dentro de la aplicación.
\end{itemize}

Con este umbral cubierto, el siguiente salto funcional razonable del proyecto no pasa ya por incorporar \texttt{M4}, sino por consolidar documentalmente su alcance final y, \textcolor{red}{en futuras iteraciones, ampliar el dominio de cobros hacia pagos parciales, facturación y conciliación avanzada si el proyecto lo requiere.}
